% !TeX root = ../../main.tex
\subsection{2乗に比例する関数について}

2つの数量 $x$，$y$ があり，$y$ が $x^2$ に比例するとき，

\[y = ax^2 \qquad (a \neq 0)\]

という関係が成り立つ.

\begin{definitionbox}{\textbf{2乗に比例する関数}}
    $y$ が $x^2$ に比例するとき，$y = ax^2$（$a \neq 0$）という関係が成り立つ.
    このグラフは\textbf{放物線}を描く.
\end{definitionbox}

\begin{theorembox}{\textbf{放物線の性質}}
    $y = ax^2$ のグラフについて，
    \begin{enumerate}
        \item $a > 0$ のとき，下に凸で原点が\textbf{最小値}をとる.
        \item $a < 0$ のとき，上に凸で原点が\textbf{最大値}をとる.
        \item $A \leqq x \leqq B$ のとき，$y$ の変域は $C \leqq y \leqq D$ となる.
        \item 原点を通る.
    \end{enumerate}
\end{theorembox}

\begin{theorembox}{\textbf{$x$ と $y$ の対応}}
    $y = ax^2$ において，$x = p$，$x = q$ のときの $y$ の値の和は $a(p^2 + q^2)$ である.
\end{theorembox}

\begin{example}[【例題 1】]
    $y = 2x^2$ において，$x = -3$ のときの $y$ の値を求めなさい.
\end{example}

\begin{proof}\mbox{}\\
    式に代入する.

    \begin{hiddenanswer}
        \begin{align*}
            y &= 2 \times (-3)^2 \\
              &= 2 \times 9 \\
              &= 18
        \end{align*}
        \finalanswer{答え：$18$}
    \end{hiddenanswer}
\end{proof}

\newpage
